%% ===================================================================================
%% Capítulo 4 -- Levantamento e Análise de Requisitos
%% FONTE DA VERDADE dos requisitos, das regras de negócio e dos casos de uso.
%% Origem: Relatório de IC (RF01 a RF04, diagrama de casos de uso, protótipos) + derivação
%% feita na consolidação do monólito (RF05 em diante, RNF e RN).
%% Dono: Arquiteto/Redator. Desenvolvedores: citem os identificadores (RF, RNF, RN, UC)
%% nos commits e nos pedidos do mural; peçam mudanças pelo mural.
%% ===================================================================================
\chapter{Levantamento e Análise de Requisitos}\label{cap:requisitos}

Este capítulo especifica \emph{o que} o Vitrine deve fazer e sob que restrições. Os requisitos RF01 a RF04 foram definidos no relatório de IC \cite{reis2026ic}; os demais requisitos funcionais, os requisitos não funcionais e as regras de negócio foram derivados dos casos de uso, dos protótipos e dos objetivos (Seção~\ref{sec:objetivos-especificos}) durante a consolidação do monólito, e estão marcados com a sua origem. Todos os requisitos e regras deste capítulo, inclusive RF26 e RN15 a RN17, foram aprovados pelo Engenheiro Chefe e são oficiais. A especificação segue as recomendações de identificação única, verificabilidade e rastreabilidade da ISO/IEC/IEEE 29148 \cite{iso29148}.

\section{Visão Geral do Produto}\label{sec:visao-produto}

\textbf{Propósito.} Permitir que docentes e pesquisadores publiquem e organizem conteúdo científico e didático --- artigos, notas de aula, tutoriais --- em um site próprio, com identidade visual personalizada, edição em Markdown com fórmulas \LaTeX{} e entrega com desempenho de página estática.

\textbf{Perspectiva.} O Vitrine é uma aplicação web autônoma, implantada como um contêiner Docker com um único processo Node.js, acompanhado de um arquivo de banco de dados SQLite em volume persistente. Não depende de serviços externos obrigatórios (Figura~\ref{fig:arq-contexto}, Capítulo~\ref{cap:arquitetura}).

\textbf{Ambiente operacional.} Navegadores modernos no cliente; no servidor, uma VPS tradicional com Docker, volume persistente em disco local e, em produção, HTTPS (decisão \decisao{D-04}, \adr{ADR-011}).

\textbf{Restrições.} (i) A pilha tecnológica é Next.js, Prisma, SQLite e better-auth; (ii) a interface é em português do Brasil; (iii) o código é aberto, sob licença MIT (\adr{ADR-015}); (iv) a capacidade de entrega do time é limitada, o que exige priorização (\emph{risco} R10, Tabela~\ref{tab:riscos}).

\textbf{Premissas e dependências.} Assume-se a disponibilidade de uma VPS com Docker e armazenamento persistente local (\decisao{D-04}); a estabilidade das interfaces de programação das versões fixadas (Tabela~\ref{tab:stack}); e a existência de um único administrador na versão inicial (\decisao{D-13}).

\section{Atores e Partes Interessadas}\label{sec:atores}

\begin{table}[htb]
\caption{Atores do sistema e partes interessadas}\label{tab:atores}
\small
\begin{tabularx}{\linewidth}{@{}L{2.8cm}L{2.4cm}Y@{}}
\toprule
\textbf{Ator ou parte} & \textbf{Tipo} & \textbf{Descrição e interesse}\\
\midrule
Leitor & Ator primário & Visitante --- estudante, colega, pesquisador --- que navega e lê o conteúdo publicado; quer encontrar e ler rápido, em qualquer dispositivo. Não se autentica.\\
\addlinespace
Administrador & Ator primário & Docente-autor que cria, organiza e publica conteúdo e define a identidade visual do site. Tem conhecimentos básicos de informática e não deseja administrar infraestrutura.\\
\addlinespace
Sistema & Ator secundário & Automação interna que dispara a invalidação do cache ISR após mutações (caso de uso UC12).\\
\addlinespace
Implantador & Parte interessada & Quem instala e opera o sistema --- em geral o próprio docente com apoio técnico ou um desenvolvedor. Interessa-se por simplicidade de instalação, configuração e \en{backup}.\\
\addlinespace
Engenheiro Chefe & Parte interessada & Autor e responsável pelo produto; decide o escopo e as questões de projeto e aprova as integrações e as regras de negócio.\\
\addlinespace
Agentes de desenvolvimento & Parte interessada & Dados, Core e UI; consomem esta especificação como contrato de implementação.\\
\bottomrule
\end{tabularx}
\fonte{Elaborado pelos autores, com base em \citeonline{reis2026ic}.}
\end{table}

\section{Casos de Uso}\label{sec:casos-uso}

\subsection{Diagrama de Casos de Uso}

A Figura~\ref{fig:uc-visao-geral} apresenta o diagrama de casos de uso do Vitrine. Ele parte do diagrama do relatório de IC (Figura~\ref{fig:ic-uc}), com uma alteração: \emph{Gerenciar Conteúdo} passa a incluir \emph{Invalidar Cache (ISR)}, porque toda mutação de conteúdo --- e não apenas a personalização --- deve refletir-se nas páginas públicas (\req{RF17}).

\vitrinefigura[tipo=diagrama, arquivo=uc-visao-geral, agente=Arquiteto, largura=0.95,
  origem={doc/diagrams/src/uc-visao-geral.puml},
  descricao={Diagrama de casos de uso com os atores Leitor, Administrador e Sistema, os pacotes Navegação e Leitura, Gestão de Estrutura e Configurações e Performance, e as relações include e extend.}]
  {Diagrama de casos de uso do Vitrine}{fig:uc-visao-geral}

\vitrinefigura[tipo=ic, arquivo=ic-fig1-casos-de-uso, largura=0.4,
  fonte={Relatório de IC \cite{reis2026ic}, Figura 1.}]
  {Diagrama de casos de uso original, do relatório de IC}{fig:ic-uc}

\subsection{Catálogo de Casos de Uso}

\begin{longtable}{@{}L{1.2cm}L{4.4cm}L{2.5cm}L{5.0cm}@{}}
\caption{Catálogo de casos de uso}\label{tab:catalogo-uc}\\
\toprule
\textbf{ID} & \textbf{Caso de uso} & \textbf{Ator} & \textbf{Descrição resumida}\\
\midrule
\endfirsthead
\toprule
\textbf{ID} & \textbf{Caso de uso} & \textbf{Ator} & \textbf{Descrição resumida}\\
\midrule
\endhead
\bottomrule
\endfoot
UC01 & Explorar \en{sidebar} hierárquica & Leitor & Navegar pela árvore de coleções e documentos e pelas coleções raiz na barra superior.\\
UC02 & Visualizar página única & Leitor & Abrir um documento por sua URL amigável.\\
UC03 & Ler conteúdo (Markdown) & Leitor & Ler o documento renderizado, com fórmulas, tabelas e código. Incluído por UC01 e UC02.\\
UC04 & Gerenciar conteúdo & Administrador & Criar, editar, definir o estado de publicação (rascunho ou publicado), mover, ordenar e excluir conteúdos no painel. Inclui UC07, UC08 e UC12.\\
UC05 & Criar documento & Administrador & Criar conteúdo do tipo DOCUMENT. Estende UC04.\\
UC06 & Criar coleção/pasta & Administrador & Criar conteúdo do tipo COLLECTION. Estende UC04.\\
UC07 & Aninhar e ordenar conteúdos & Administrador & Definir pai e ordem de um conteúdo. Incluído por UC04.\\
UC08 & Autenticar-se & Administrador & Entrar e sair da área administrativa. Incluído por UC04 e UC09.\\
UC09 & Personalizar UI & Administrador & Alterar a aparência do site. Inclui UC08 e UC12.\\
UC10 & Definir identidade (nome e \en{branding}) & Administrador & Definir nome, descrição, logotipo, \en{favicon} e links sociais do rodapé. Estende UC09.\\
UC11 & Configurar paleta de cores & Administrador & Escolher as cores do tema. Estende UC09.\\
UC12 & Invalidar cache (ISR) & Sistema & Invalidar as páginas afetadas após uma mutação.\\
\end{longtable}
\fonte{Elaborado pelos autores, a partir da Figura~\ref{fig:uc-visao-geral}.}

\subsection{Especificações Detalhadas}\label{sec:uc-detalhados}

Os casos de uso com maior impacto sobre a arquitetura são especificados a seguir. Os demais seguem os mesmos padrões e são detalhados conforme a necessidade de cada módulo.

\begin{table}[htb]
\caption{Especificação do caso de uso UC08 --- Autenticar-se}\label{tab:uc08}
\small
\begin{tabularx}{\linewidth}{@{}L{3.0cm}Y@{}}
\toprule
\textbf{Ator} & Administrador\\
\textbf{Pré-condição} & Existe ao menos um usuário cadastrado (bootstrap, \req{RF22}).\\
\textbf{Gatilho} & Acesso a uma rota protegida sem sessão válida, ou acesso direto a \arq{/login}.\\
\textbf{Fluxo principal} & (1) O sistema exibe o formulário de e-mail e senha. (2) O administrador informa as credenciais. (3) O sistema verifica as credenciais e cria uma sessão persistida no banco. (4) O sistema grava o \en{cookie} de sessão (\en{HttpOnly}) e redireciona ao destino original ou a \arq{/admin}.\\
\textbf{Fluxos alternativos} & A1: credenciais inválidas --- mensagem genérica, que não revela qual campo falhou; limitação de tentativas (\req{RNF03}). A2: sessão expirada --- retorno ao formulário preservando o destino (parâmetro \texttt{redirectTo}, que aceita apenas caminhos internos do site; destinos externos levam ao painel). A3: encerrar sessão (\req{RF20}) --- a sessão é invalidada e o \en{cookie} removido.\\
\textbf{Pós-condição} & Sessão válida associada ao usuário; operações de escrita autorizadas (\req{RF21}).\\
\textbf{Requisitos e regras} & \req{RF01}, \req{RF20}, \req{RF21}, \req{RF22}; RN01, RN12.\\
\bottomrule
\end{tabularx}
\fonte{Elaborado pelos autores.}
\end{table}

\begin{table}[htb]
\caption{Especificação do caso de uso UC05 --- Criar documento}\label{tab:uc05}
\small
\begin{tabularx}{\linewidth}{@{}L{3.0cm}Y@{}}
\toprule
\textbf{Ator} & Administrador (autenticado, UC08)\\
\textbf{Pré-condição} & Sessão válida. O documento será criado na raiz ou dentro de uma coleção existente.\\
\textbf{Fluxo principal} & (1) O administrador aciona ``Nova publicação'' (na raiz ou na coleção desejada). (2) Informa o título; o \en{slug} é sugerido a partir do título e pode ser editado; escolhe a coleção pai e o estado de publicação (padrão: Publicado). (3) Escreve o texto em Markdown, acompanhando a pré-visualização. (4) Aciona ``Enviar nova versão''. (5) O sistema valida os dados (RN02, RN03, RN06), persiste o conteúdo com o usuário como autor e (6) dispara a invalidação do cache (UC12). (7) O painel confirma a publicação.\\
\textbf{Fluxos alternativos} & A1: \en{slug} já existente entre os irmãos --- o sistema responde com conflito e sugere um \en{slug} alternativo. A2: o pai escolhido não é uma coleção --- erro de validação. A3: o administrador opta por ``Salvar localmente'' e retoma depois (\req{RF15}).\\
\textbf{Pós-condição} & Documento persistido. Se ele e todos os seus ancestrais estiverem publicados, a \en{sidebar} e a página refletem a mudança a partir da próxima requisição; um rascunho só é visível no painel (RN15).\\
\textbf{Requisitos e regras} & \req{RF05}, \req{RF08}, \req{RF09}, \req{RF15}, \req{RF16}, \req{RF17}, \req{RF26}; RN01 a RN06, RN10, RN11, RN13, RN15, RN16.\\
\bottomrule
\end{tabularx}
\fonte{Elaborado pelos autores.}
\end{table}

\begin{table}[htb]
\caption{Especificação do caso de uso UC07 --- Aninhar e ordenar conteúdos}\label{tab:uc07}
\small
\begin{tabularx}{\linewidth}{@{}L{3.0cm}Y@{}}
\toprule
\textbf{Ator} & Administrador (autenticado, UC08)\\
\textbf{Pré-condição} & Sessão válida; existem conteúdos a reorganizar.\\
\textbf{Fluxo principal} & (1) O administrador visualiza a árvore de conteúdos no painel. (2) Arrasta um item para outra posição ou outro pai, ou usa as ações ``subir'' e ``descer''. (3) A interface atualiza a árvore de forma otimista e envia o lote de alterações. (4) O sistema valida o lote (pai do tipo coleção, ausência de ciclos, \en{slugs} únicos no destino) e o persiste em uma transação. (5) O sistema invalida o \en{layout} raiz (UC12) e devolve a árvore atualizada.\\
\textbf{Fluxos alternativos} & A1: movimento inválido --- o sistema rejeita o lote inteiro, a interface desfaz a atualização otimista e exibe o motivo.\\
\textbf{Pós-condição} & Nova estrutura e nova ordem persistidas; a \en{sidebar} pública reflete a mudança.\\
\textbf{Requisitos e regras} & \req{RF02}, \req{RF03}, \req{RF17}; RN02, RN03, RN05, RN06.\\
\bottomrule
\end{tabularx}
\fonte{Elaborado pelos autores.}
\end{table}

\begin{table}[htb]
\caption{Especificação do caso de uso UC11 --- Configurar paleta de cores}\label{tab:uc11}
\small
\begin{tabularx}{\linewidth}{@{}L{3.0cm}Y@{}}
\toprule
\textbf{Ator} & Administrador (autenticado, UC08)\\
\textbf{Pré-condição} & Sessão válida.\\
\textbf{Fluxo principal} & (1) O administrador abre Configurações. (2) Escolhe as cores primária, secundária e de destaque por seletor ou por código hexadecimal. (3) Pré-visualiza o resultado no próprio painel, com variáveis CSS locais. (4) Salva. (5) O sistema valida o formato e calcula o contraste, devolvendo avisos se for o caso. (6) Persiste a configuração e invalida o \en{layout} raiz. (7) Os leitores veem a nova paleta na próxima requisição, sem novo \en{build}.\\
\textbf{Fluxos alternativos} & A1: valor fora do formato \texttt{\#RRGGBB} --- erro de validação no campo. A2: contraste abaixo de 4,5:1 entre texto e fundo --- aviso, sem bloquear o salvamento (RN09).\\
\textbf{Pós-condição} & Configuração persistida; páginas regeneradas com as novas variáveis CSS.\\
\textbf{Requisitos e regras} & \req{RF04}, \req{RF17}, \req{RF19}; RN08, RN09, RN10.\\
\bottomrule
\end{tabularx}
\fonte{Elaborado pelos autores.}
\end{table}

\section{Requisitos Funcionais}\label{sec:rf}

A Tabela~\ref{tab:rf} lista os requisitos funcionais. A coluna \emph{Origem} indica se o requisito foi herdado do relatório de IC (``IC'') ou derivado nesta especificação a partir de casos de uso (UC), protótipos, objetivos (OE), regras (RN) ou decisões (D). A prioridade segue a escala \emph{Essencial}, \emph{Importante} e \emph{Desejável}. A coluna \emph{Agente} indica o responsável principal pela implementação; requisitos que atravessam domínios são coordenados pelo mural.

\begin{longtable}{@{}L{1.3cm}L{7.2cm}L{1.6cm}L{2.2cm}L{1.8cm}@{}}
\caption{Requisitos funcionais}\label{tab:rf}\\
\toprule
\textbf{ID} & \textbf{Requisito} & \textbf{Prior.} & \textbf{Origem} & \textbf{Agente}\\
\midrule
\endfirsthead
\toprule
\textbf{ID} & \textbf{Requisito} & \textbf{Prior.} & \textbf{Origem} & \textbf{Agente}\\
\midrule
\endhead
\bottomrule
\endfoot
\reqrow{RF01}{Autenticação}{O acesso à área administrativa exige e-mail e senha, com sessão persistida em banco (better-auth).}{Essencial}{IC (RF01); UC08}{Core}
\reqrow{RF02}{Estrutura hierárquica}{Os conteúdos formam uma árvore de coleções e documentos aninhados, por auto-relacionamento (pai e filhos).}{Essencial}{IC (RF02); UC04, UC07}{Dados}
\reqrow{RF03}{Ordenação}{Cada conteúdo possui uma ordem relativa entre irmãos, editável pelo administrador.}{Essencial}{IC (RF03); UC07}{Dados}
\reqrow{RF04}{Personalização}{Interface para alterar identidade visual e cores por variáveis CSS. Requisito agregador, detalhado em RF18 e RF19.}{Essencial}{IC (RF04); UC09}{UI}
\reqrow{RF05}{Criar documento}{Criar conteúdo do tipo DOCUMENT com título, \en{slug}, pai, estado de publicação (padrão PUBLISHED) e corpo em Markdown.}{Essencial}{UC05}{Core}
\reqrow{RF06}{Criar coleção}{Criar conteúdo do tipo COLLECTION (pasta) para agrupar documentos e subcoleções, com estado de publicação (padrão PUBLISHED).}{Essencial}{UC06}{Core}
\reqrow{RF07}{Editar e excluir}{Atualizar título, \en{slug}, corpo, estado, pai e ordem; excluir com confirmação e remoção em cascata dos descendentes.}{Essencial}{UC04; RN07}{Core}
\reqrow{RF08}{Editor Markdown}{Editor com barra de ferramentas e pré-visualização em tempo real (painel dividido).}{Essencial}{IC (OE2); Fig. 6}{UI}
\reqrow{RF09}{Fórmulas matemáticas}{Renderizar fórmulas \LaTeX{} em linha e em bloco no conteúdo e na pré-visualização (KaTeX).}{Essencial}{IC (OE2); UC03}{UI}
\reqrow{RF10}{Realce de sintaxe}{Realçar blocos de código por linguagem, no conteúdo e na pré-visualização.}{Importante}{IC (OE2); UC03}{UI}
\reqrow{RF11}{Markdown estendido}{Suportar tabelas, listas de tarefas, texto tachado e links automáticos (GFM).}{Importante}{IC (Fig. 3); UC03}{UI}
\reqrow{RF12}{\en{Sidebar} hierárquica}{Exibir a árvore da coleção corrente, com o item ativo destacado, níveis recolhíveis e navegação por teclado.}{Essencial}{UC01; Fig. 3}{UI}
\reqrow{RF13}{Barra superior}{Exibir logotipo, nome do site e as coleções raiz como abas ou links de navegação.}{Essencial}{UC01; Fig. 3}{UI}
\reqrow{RF14}{Página de documento}{Exibir o documento em URL derivada da cadeia de \en{slugs}; responder 404 a caminhos inexistentes ou não publicados (RN15).}{Essencial}{UC02; RN04}{Core}
\reqrow{RF15}{Rascunho local}{Guardar rascunho no navegador (``Salvar localmente'') sem publicá-lo.}{Desejável}{IC (Fig. 6)}{UI}
\reqrow{RF16}{Publicar versão}{Enviar ao servidor a versão editada (``Enviar nova versão''), substituindo o conteúdo publicado.}{Essencial}{IC (Fig. 6); UC05}{Core}
\reqrow{RF17}{Invalidação de cache}{Revalidar as páginas afetadas após qualquer mutação de conteúdo ou de configuração.}{Essencial}{IC (UC12); RN10}{Core}
\reqrow{RF18}{Identidade do site}{Configurar nome, descrição, logotipo e \en{favicon}.}{Essencial}{UC10; Fig. 5}{Core}
\reqrow{RF19}{Paleta de cores}{Configurar as cores primária, secundária e de destaque, aplicadas por variáveis CSS sem novo \en{build}.}{Essencial}{IC (OE3); UC11}{UI}
\reqrow{RF20}{Encerrar sessão}{Permitir \en{logout}, invalidando a sessão no servidor.}{Essencial}{UC08}{Core}
\reqrow{RF21}{Autorização de escrita}{Bloquear, no servidor, toda operação de escrita sem sessão válida.}{Essencial}{UC08; RN01}{Core}
\reqrow{RF22}{Bootstrap do administrador}{Criar o primeiro administrador na instalação, pelo comando \texttt{npm run admin:create}, a partir das variáveis de ambiente \texttt{ADMIN\_EMAIL}, \texttt{ADMIN\_PASSWORD} e, opcionalmente, \texttt{ADMIN\_NAME}; o comando é idempotente, recusa um segundo usuário com outro e-mail e não imprime a senha. O cadastro público é desabilitado e não há página de primeiro acesso.}{Essencial}{D-03; ADR-018}{Core}
\reqrow{RF23}{Configuração inicial}{Inserir a configuração padrão (registro único) na instalação.}{Essencial}{RN08}{Dados}
\reqrow{RF24}{Rodapé configurável}{Exibir rodapé com os links sociais configurados: seis campos opcionais e independentes (Discord, GitHub, YouTube, Instagram, LinkedIn e Twitch), apenas URLs http(s); campo vazio não exibe o ícone (RN17). O texto de direitos é gerado a partir do nome do site e do ano.}{Desejável}{IC (Fig. 3 a 5); D-19; ADR-013}{UI}
\reqrow{RF25}{Metadados de SEO}{Gerar título, descrição e metadados por página; expor \arq{robots.txt} e mapa do site, este somente com conteúdo publicado (RN15).}{Importante}{IC (OE1); RNF07}{Core}
\reqrow{RF26}{Estado de publicação}{Cada conteúdo tem estado DRAFT (rascunho) ou PUBLISHED (publicado), com padrão PUBLISHED. Somente conteúdo publicado, com todos os ancestrais publicados, é exibido ao público; o painel mostra todos os conteúdos e permite alternar o estado.}{Essencial}{D-15; ADR-012; RN15, RN16}{Core}
\end{longtable}
\fonte{Elaborado pelos autores. RF01 a RF04: \citeonline{reis2026ic}.}

\section{Requisitos Não Funcionais}\label{sec:rnf}

Os requisitos não funcionais (Tabela~\ref{tab:rnf}) estão associados às características de qualidade da ISO/IEC 25010 \cite{iso25010}. As metas numéricas são as metas oficiais do produto e são verificadas no Capítulo~\ref{cap:testes}.

\begin{longtable}{@{}L{1.4cm}L{7.6cm}L{2.7cm}L{1.9cm}@{}}
\caption{Requisitos não funcionais}\label{tab:rnf}\\
\toprule
\textbf{ID} & \textbf{Requisito} & \textbf{Característica (ISO 25010)} & \textbf{Prior.}\\
\midrule
\endfirsthead
\toprule
\textbf{ID} & \textbf{Requisito} & \textbf{Característica (ISO 25010)} & \textbf{Prior.}\\
\midrule
\endhead
\bottomrule
\endfoot
\reqtarget{RNF01} & \textbf{Portabilidade e implantação simples.} Implantável a partir do repositório com poucos comandos (instalar, configurar variáveis de ambiente, migrar, compilar e iniciar), com um único arquivo de banco e nenhum serviço externo obrigatório. & Portabilidade & Essencial\\[1pt]
\reqtarget{RNF02} & \textbf{Desempenho de leitura.} Páginas públicas servidas do cache ISR. Meta: LCP $\leq$ 2,5 s (percentil 75) e pontuação de desempenho $\geq$ 90 no Lighthouse, em perfil móvel simulado. & Eficiência de desempenho & Importante\\[1pt]
\reqtarget{RNF03} & \textbf{Segurança.} Senhas somente como \en{hash}; \en{cookies} de sessão \texttt{HttpOnly}, \texttt{SameSite} e \texttt{Secure} em produção; validação de toda entrada no servidor; proibição de HTML bruto no conteúdo; limitação de tentativas de login (no máximo cinco por minuto por endereço IP); recusa de mutações vindas de outra origem; segredos apenas em variáveis de ambiente; cabeçalhos de segurança HTTP; alinhamento ao OWASP Top 10 \cite{owasptop10}. & Segurança & Essencial\\[1pt]
\reqtarget{RNF04} & \textbf{Usabilidade.} Interface em português do Brasil; fluxo de criar e publicar um documento concluível por usuário sem formação técnica e sem consultar documentação. Meta: pontuação SUS $\geq$ 68. & Usabilidade & Importante\\[1pt]
\reqtarget{RNF05} & \textbf{Acessibilidade.} Conformidade com o WCAG 2.1 nível AA \cite{wcag21} nas páginas públicas e no painel: contraste mínimo de 4,5:1, navegação por teclado, foco visível, rótulos e textos alternativos. & Usabilidade & Importante\\[1pt]
\reqtarget{RNF06} & \textbf{Responsividade.} Layout adaptável de 320 px a telas largas; \en{sidebar} recolhível em telas pequenas. & Usabilidade; Compatibilidade & Importante\\[1pt]
\reqtarget{RNF07} & \textbf{SEO.} HTML pré-renderizado com título, descrição e metadados por página; URLs estáveis e legíveis. & Adequação funcional & Importante\\[1pt]
\reqtarget{RNF08} & \textbf{Manutenibilidade e modularidade.} Fronteiras entre Dados, Core e UI respeitadas; acesso ao banco somente por \arq{src/lib/db.js}; \en{lint} sem erros; documentação sincronizada com o código. & Manutenibilidade & Essencial\\[1pt]
\reqtarget{RNF09} & \textbf{Extensibilidade.} Novos \en{plugins} \texttt{remark}/\texttt{rehype}, tipos de conteúdo e campos de configuração adicionáveis sem alterar a estrutura central. & Manutenibilidade & Desejável\\[1pt]
\reqtarget{RNF10} & \textbf{Integridade e confiabilidade.} Operações compostas (mover, ordenar) em transação; rotina documentada de \en{backup} e restauração; falha na revalidação não pode causar perda de dados. & Confiabilidade & Importante\\[1pt]
\reqtarget{RNF11} & \textbf{Compatibilidade de navegadores.} Duas últimas versões estáveis de Chrome, Firefox, Safari e Edge. & Compatibilidade & Importante\\[1pt]
\reqtarget{RNF12} & \textbf{Internacionalização.} Nesta versão apenas pt-BR; textos de interface centralizados para viabilizar tradução futura. & Adequação funcional & Desejável\\[1pt]
\reqtarget{RNF13} & \textbf{Licenciamento.} Código publicado sob licença de código aberto (D-02), com dependências de licenças compatíveis. & Restrição & Essencial\\[1pt]
\reqtarget{RNF14} & \textbf{Observabilidade mínima.} Erros de servidor registrados em \en{log} com contexto suficiente ao diagnóstico, sem expor dados sensíveis ao usuário. & Manutenibilidade & Desejável\\[1pt]
\reqtarget{RNF15} & \textbf{Privacidade.} Não coletar dados pessoais de leitores além dos registros de acesso do servidor; sem rastreadores de terceiros por padrão; dados do administrador limitados a nome e e-mail. & Segurança & Importante\\[1pt]
\end{longtable}
\fonte{Elaborado pelos autores, com base em \citeonline{iso25010}.}

\section{Regras de Negócio}\label{sec:regras-negocio}

As regras de negócio (RN) constituem invariantes do domínio; devem ser garantidas no servidor, independentemente da interface.

\begin{longtable}{@{}L{1.3cm}L{9.6cm}L{2.8cm}@{}}
\caption{Regras de negócio}\label{tab:rn}\\
\toprule
\textbf{ID} & \textbf{Regra} & \textbf{Relacionados}\\
\midrule
\endfirsthead
\toprule
\textbf{ID} & \textbf{Regra} & \textbf{Relacionados}\\
\midrule
\endhead
\bottomrule
\endfoot
\reqtarget{RN01} & Somente usuários autenticados podem criar, alterar, mover ou excluir conteúdo e alterar a configuração. & RF21; UC08\\[1pt]
\reqtarget{RN02} & Todo conteúdo tem tipo DOCUMENT ou COLLECTION. Apenas COLLECTION pode ter filhos; DOCUMENT é folha. O tipo é validado na aplicação, pois o SQLite não impõe valores de enumeração. & RF02, RF05, RF06\\[1pt]
\reqtarget{RN03} & O \en{slug} contém apenas letras minúsculas, dígitos e hífens (\texttt{\^{}[a-z0-9]+(-[a-z0-9]+)*\$}), com 1 a 80 caracteres; é único entre irmãos; e os \en{slugs} \texttt{admin}, \texttt{api}, \texttt{login} e \arq{_next} são reservados no nível raiz (D-05). & RF05, RF14\\[1pt]
\reqtarget{RN04} & A URL pública de um conteúdo é a concatenação dos \en{slugs} da raiz até o nó, separados por barras. & RF14\\[1pt]
\reqtarget{RN05} & A ordem é um inteiro não negativo; entre irmãos, a listagem ordena por ordem crescente e, em empate, por data de criação. & RF03\\[1pt]
\reqtarget{RN06} & Mover um conteúdo não pode criar ciclo: o destino não pode ser o próprio nó nem um de seus descendentes. & RF02, RF03\\[1pt]
\reqtarget{RN07} & Excluir um conteúdo exclui todos os seus descendentes e exige confirmação explícita no painel. & RF07\\[1pt]
\reqtarget{RN08} & Existe exatamente um registro de configuração; se ausente, é criado com os valores padrão (nome ``Vitrine'', cores \texttt{\#3b82f6}, \texttt{\#1e293b} e \texttt{\#f59e0b}). & RF18, RF19, RF23\\[1pt]
\reqtarget{RN09} & Cores são aceitas no formato \texttt{\#RRGGBB}; o sistema avisa, sem bloquear, quando o contraste de uma cor do tema com o branco é inferior a 4,5:1 (primária e secundária, usadas com texto) ou a 3:1 (destaque, que não é texto); os pares estão na Tabela~\ref{tab:pares-contraste}. & RF19; RNF05\\[1pt]
\reqtarget{RN10} & Toda mutação de conteúdo ou de configuração dispara a revalidação das páginas afetadas: somente a página, se apenas o corpo mudou; o \en{layout} raiz, se título, \en{slug}, pai, ordem, estado de publicação, tipo (documento ou coleção) ou configuração mudaram. & RF17\\[1pt]
\reqtarget{RN11} & O corpo em Markdown nunca é interpretado como HTML bruto. & RNF03\\[1pt]
\reqtarget{RN12} & Sessões têm prazo de expiração: sete dias, renovados no máximo uma vez por dia enquanto a sessão é usada; uma sessão expirada equivale a nenhuma sessão. & RF01\\[1pt]
\reqtarget{RN13} & O autor de um conteúdo é o usuário autenticado que o criou. & RF05\\[1pt]
\reqtarget{RN14} & A abertura da URL de uma COLLECTION publicada redireciona ao primeiro documento descendente publicado, na ordem da árvore; se não houver, exibe um estado vazio (D-16). & RF14\\[1pt]
\reqtarget{RN15} & Um conteúdo é visível ao público --- por URL, na \en{sidebar} e na barra superior, na geração estática e no mapa do site --- somente se ele e todos os seus ancestrais estiverem PUBLISHED. Conteúdo não visível responde 404 por URL direta. O painel administrativo exibe todos os conteúdos, com o seu estado. & RF14, RF25, RF26\\[1pt]
\reqtarget{RN16} & O estado padrão na criação é PUBLISHED. O administrador pode alternar entre DRAFT e PUBLISHED a qualquer momento. Alterar o estado de uma COLLECTION afeta a visibilidade de toda a sua subárvore, sem modificar o estado dos descendentes. Mudar o estado é uma mudança estrutural para a revalidação (RN10). & RF26, RF17\\[1pt]
\reqtarget{RN17} & Os links sociais do rodapé (Discord, GitHub, YouTube, Instagram, LinkedIn e Twitch) são opcionais e independentes. Só são aceitas URLs absolutas com protocolo \texttt{http} ou \texttt{https}; campo vazio equivale a link ausente e o ícone não é exibido; os links são renderizados com os atributos \texttt{noopener} e \texttt{noreferrer}. & RF24, RNF03\\[1pt]
\end{longtable}
\fonte{Elaborado pelos autores.}

\vitrinependencia{Dados}{A regra RN02 (apenas COLLECTION tem filhos) e a unicidade de slug entre irmãos (RN03) ainda não estão garantidas pelo esquema atual; decidir entre restrição no banco e validação na aplicação (decisões D-05 e D-14) e registrar no mural.}

\section{Matriz de Rastreabilidade}\label{sec:matriz-rastreabilidade}

A Tabela~\ref{tab:matriz-rastreabilidade} relaciona cada requisito funcional ao caso de uso e aos módulos de implementação (definidos na Seção~\ref{sec:modulos}). A relação com os casos de teste está no Capítulo~\ref{cap:testes}.

\begin{longtable}{@{}L{1.4cm}L{2.5cm}L{3.8cm}L{3.0cm}@{}}
\caption{Rastreabilidade entre requisitos funcionais, casos de uso e módulos}\label{tab:matriz-rastreabilidade}\\
\toprule
\textbf{RF} & \textbf{Caso(s) de uso} & \textbf{Módulo(s)} & \textbf{Agente(s)}\\
\midrule
\endfirsthead
\toprule
\textbf{RF} & \textbf{Caso(s) de uso} & \textbf{Módulo(s)} & \textbf{Agente(s)}\\
\midrule
\endhead
\bottomrule
\endfoot
RF01 & UC08 & M03 & Core\\
RF02 & UC04, UC07 & M01, M02 & Dados\\
RF03 & UC07 & M01, M02, M10 & Dados, UI\\
RF04 & UC09 & M11 & UI, Core\\
RF05 & UC05 & M04, M09 & Core, UI\\
RF06 & UC06 & M04, M10 & Core, UI\\
RF07 & UC04 & M04, M10 & Core, UI\\
RF08 & UC04 & M09 & UI\\
RF09 & UC03 & M08 & UI\\
RF10 & UC03 & M08 & UI\\
RF11 & UC03 & M08 & UI\\
RF12 & UC01 & M07 & UI\\
RF13 & UC01 & M07 & UI\\
RF14 & UC02 & M02, M07 & Core, Dados\\
RF15 & UC05 & M09 & UI\\
RF16 & UC05 & M04, M09 & Core, UI\\
RF17 & UC12 & M06 & Core\\
RF18 & UC10 & M05, M11 & Core, UI\\
RF19 & UC11 & M05, M11 & UI, Core\\
RF20 & UC08 & M03 & Core\\
RF21 & UC08 & M03, M04, M05 & Core\\
RF22 & UC08 & M01, M03, M13 & Core, Dados\\
RF23 & UC09 & M01 & Dados\\
RF24 & UC01, UC10 & M01, M07, M11 & UI, Dados\\
RF25 & UC02 & M12 & Core\\
RF26 & UC02, UC04, UC05 & M01, M04, M07, M10 & Core, Dados, UI\\
\end{longtable}
\fonte{Elaborado pelos autores.}

\section{Prototipagem de Baixa Fidelidade}\label{sec:prototipagem}

A prototipagem de baixa fidelidade permitiu validar a jornada do usuário e a disposição dos elementos antes da codificação final \cite{reis2026ic}. As quatro interfaces do relatório de IC, reproduzidas a seguir, são a referência visual para a implementação da interface (agente UI) e originaram requisitos funcionais.

\subsection{Visualização do Sistema com Conteúdo}

A Figura~\ref{fig:ic-instancia} mostra uma instância do sistema em operação, preenchida com conteúdos de programação. Ela evidencia a barra superior com as coleções raiz (RF13), a \en{sidebar} hierárquica com níveis recolhíveis (RF12), a renderização de títulos, listas, tabelas e texto tachado (RF11) e o rodapé com ícones de redes sociais (RF24).

\vitrinefigura[tipo=ic, arquivo=ic-fig3-instancia-fromcaio, largura=0.7,
  fonte={Relatório de IC \cite{reis2026ic}, Figura 3.}]
  {Instância do sistema populada com cursos de programação (protótipo da IC)}{fig:ic-instancia}

\subsection{Gestão e Organização de Conteúdos}

O sistema permite uma estrutura hierárquica em que coleções e arquivos são aninhados. A Figura~\ref{fig:ic-gestao} apresenta a interface administrativa de organização, em que o usuário gerencia a árvore e define a precedência de exibição de cada item, com ações de criar, editar e excluir por nó (RF03, RF06 e RF07).

\vitrinefigura[tipo=ic, arquivo=ic-fig4-gestao-conteudo, largura=0.6,
  fonte={Relatório de IC \cite{reis2026ic}, Figura 4.}]
  {Interface de organização de documentos e coleções aninhadas (protótipo da IC)}{fig:ic-gestao}

\subsection{Customização Dinâmica}

A tela de customização (Figura~\ref{fig:ic-personalizacao}) permite ao administrador alterar a identidade visual, incluindo o logotipo e as cores. No protótipo, as cores foram rotuladas por função na barra de navegação (``principal \en{background}'' e ``principal \en{hover}''); o modelo de dados adota três cores semânticas --- primária, secundária e de destaque --- e o mapeamento entre ambos é definido na Seção~\ref{sec:arq-tema}.

\vitrinefigura[tipo=ic, arquivo=ic-fig5-personalizacao, largura=0.55,
  fonte={Relatório de IC \cite{reis2026ic}, Figura 5.}]
  {Interface de personalização de cores e identidade visual (protótipo da IC)}{fig:ic-personalizacao}

\subsection{Editor de Escrita Científica}

A interface do editor (Figura~\ref{fig:ic-editor}) foi priorizada para suportar a escrita acadêmica: Markdown integrado a fórmulas \LaTeX{} e blocos de código com realce, com barra de ferramentas, painel dividido com pré-visualização e dois comandos de publicação --- ``Salvar localmente'' e ``Enviar nova versão'' (RF08, RF15 e RF16).

\vitrinefigura[tipo=ic, arquivo=ic-fig6-editor, largura=0.65,
  fonte={Relatório de IC \cite{reis2026ic}, Figura 6.}]
  {Interface do editor de texto com suporte a elementos técnicos (protótipo da IC)}{fig:ic-editor}

\subsection{Lacunas Identificadas nos Protótipos}

A análise dos protótipos em confronto com o modelo de dados revelou duas lacunas, ambas resolvidas por decisão do Engenheiro Chefe em 30/09/2026: (i) o rodapé apresenta links de redes sociais, mas o esquema original não previa onde armazená-los, o que foi resolvido com seis campos opcionais e independentes em \texttt{Configuration} (\decisao{D-19}, \adr{ADR-013}); e (ii) o botão ``Enviar nova versão'' sugeria um histórico de versões, que o modelo de dados não contempla; adotou-se o estado de publicação \texttt{DRAFT}/\texttt{PUBLISHED}, com padrão \texttt{PUBLISHED} (\decisao{D-15}, \adr{ADR-012}), e o histórico de versões permanece fora do escopo desta versão.

\section{Síntese do Capítulo}

Foram especificados doze casos de uso, vinte e seis requisitos funcionais, quinze requisitos não funcionais e dezessete regras de negócio, com rastreabilidade para módulos e agentes. Os requisitos herdados da IC foram preservados e refinados; os derivados foram marcados. Este capítulo é o contrato de comportamento: o Capítulo~\ref{cap:arquitetura} define \emph{como} o sistema atende a ele.
